\chapter{CONCEPTS OF SNORT}
\label{chap:concepts_of_snort}

This chapter gives a brief introduction to the basic concepts of IDS's in
general and \snort{} in particular.

\section{Rules}
%TODO: Explain planning of rule modifiers and hierarchy model in Design chapter.

A \snort{} rule is made up of a header section followed by rule options. The
header specifies which network, protocols and ports the rule is valid for, and
which action to take if the rule matches traffic, while rule options provide
metadata for the rule, such as its ID and classification \cite{SnortManual}. The
rule options are specified as plain text in the {\em rule string} enclosed in
parentheses (Listing \ref{lst:rulestring}). Note that a badly formatted
rulestring given to \snort{} will cause a crash.


\begin{lstlisting}[frame=single, caption=Example rule string.,
label=lst:rulestring,captionpos=b]
alert udp $EXTERNAL_NET any -> $HOME_NET 69 \
(msg:"PROTOCOL-TFTP Get"; classtype:bad-unknown; \
sid:1444; rev:9;)
\end{lstlisting}
There are several types of rule options, but only the general rule options are
relevant for this project. The general rule options provide information about
the rule \cite{SnortManual} (i.e. signature ID, revision number and classtype).
Other types of options, such as payload options, which are actively used in
detection, are not relevant for the management of rules and are thus not considered in this
project. Some of the general options can only be specified within the
rulestring, referred to as {\em inline} specification, while others can be
specified in separate files, referred to as {\em external} specification.

A rule can either be enabled or disabled in the detection system, meaning that
the rule will only have an effect on detection when it is enabled. Since this
project is dealing with multiple sensors connected to one central database, this
functionality has been expanded so that users can keep a rule enabled on one
sensor and disabled on another. In addition, the user can also set an
enabled/disabled flag on \textit{rulesets}.

% This is our old solution:
% A rule \em modifier \em is created for each rule in a ruleset when the ruleset
% is added to a sensor. The modifier dictates whether the rule will be enabled or
% not on the given sensor, and by default it will assume the original value of
% the rule. This means that the original value is only used to determine the
% default value of the modifier. A modifier can also be configured to ``follow
% the ruleset'', in which case the rule will be enabled only when the ruleset is
% enabled, and vice versa. See table \ref{tab:ruleEnabled} for the effects of the
% modifier and ruleset options.

%Set row spacing:
% \renewcommand{\arraystretch}{1.5}
% 
% \begin{table}[tbp]
%   \centering
%   \begin{tabular}{l|l|l}
%   	Modifier & Ruleset & Active on sensor? \\
%     \hline
% 	Enabled & \em any \em & Yes \\
% 	Disabled & \em any \em & No \\
% 	Ruleset & Enabled & Yes \\
% 	Ruleset & Disabled & No \\
%   \end{tabular}
%   \caption{Determining whether a rule is enabled or disabled on a sensor.}
%   \label{tab:ruleEnabled}
% \end{table}
\section{Rulesets}

The term ruleset, sometimes rule set, is not formally defined in the
\snortManual, however it is used a few times to refer to the complete collection
of official \snort{} rules\footnote{Also referred to as the Sourcefire VRT
Certified Rules, see http://www.snort.org/snort-rules} and to the complete set
of rules applied on a \snort{} system. In practice, rules are categorized by
intrusion type (e.g. malware, trojan and worm) in order to keep the rules
structured. These groups of rules are often named by their intrusion type, and
are also referred to as rulesets. \snowman{} and this report use only the latter
definition of rulesets. The official \snort{} rules are split into files whose
names denote the ruleset. 

\section{Generators}

\textit{Generator} is a term used for any subsystem of \snort{} which generates
alerts from rules \cite{SnortManual}. Each generator is identified by a unique
generator ID. The rules subsystem, which has generator id 1, is the only
relevant generator for this project.

\section{References}
\snort{} contains a plugin which allows rules to contain
references to external attack identification systems, which provide
additional information about the alerts \cite{SnortManual}. A reference is
basically a URL string pointing to a webpage containing this information.

\section{Event and detection filters}
Filtering can be applied to rules in order to control the amount and frequency
of alerts. Two types of filters can be applied to a rule: event filters and
detection filters. A detection filter can only appear inline, whereas an event
filter can only appear as a standalone option\footnote{An event filter can
currently be written inline (called a rule threshold), but according to the
\snortManual{} this feature is deprecated and will not be supported in future
releases.}. All filters operate with a \textit{count} and \textit{seconds}
variable, representing the number of events in a time interval and the duration
of the time interval, respectively. Event filters are further divided into
three types (table \ref{tab:filterTypes}).

Detectionfilters are filters limiting the number of times a rule generates an
event, while eventfilters limits how many events are producing alarms.

\section{Rule suppression}
Some rules might match traffic that is unwanted from most hosts, while some
hosts would be allowed to send traffic matching the same rules. In this case,
\snort{} gives the possibility to remove (ie: suppress) alarms when the source
or destination of the traffic matches the configured parametres.

%Set row spacing:
\renewcommand{\arraystretch}{1.5}

\begin{table}[htb]
  \centering
  \begin{tabular}{l|l}
  	Filter type & Alerts \\
    \hline
	\textbf{Event filter} \\
	Limit & One for every event when number of events \(\leq\) C in interval S\\
	Threshold & One for every C events in interval S \\
	Both & One per interval S when number of events \(=\) C\\
	\hline
	\textbf{Detection filter} & One for every event when \textit{count} \(>\) C in
	interval S\\
	\hline
	\multicolumn{2}{l}{\textit{Note: C and S represent the \em count \em and \em
	seconds \em values in the filter.}}\\
  \end{tabular}
  \caption{Effects of different filter types.}
  \label{tab:filterTypes}
\end{table}